%!TeX root=../gardentop.tex
\chapter{A Tantrum} 
	
\begin{figure}[t!]
\centering
\includegraphics[width=\linewidth]{july}
\end{figure}

 \lettrine[]{S}{he} had got up very early in the morning and had worked hard in the garden and she was tired and sleepy, so as soon as Martha had brought her supper and she had eaten it, she was glad to go to bed. As she laid her head on the pillow she murmured to herself:

\zz
<I'll go out before breakfast and work with Dickon and then afterward—I believe—I'll go to see him.>

She thought it was the middle of the night when she was wakened by such dreadful sounds that she jumped out of bed in an instant. What was it—what was it? The next minute she felt quite sure she knew. Doors were opened and shut and there were hurrying feet in the corridors and some one was crying and screaming at the same time, screaming and crying in a horrible way.

<It's Colin,> she said. <He's having one of those tantrums the nurse called hysterics. How awful it sounds.>

As she listened to the sobbing screams she did not wonder that people were so frightened that they gave him his own way in everything rather than hear them. She put her hands over her ears and felt sick and shivering.

<I don't know what to do. I don't know what to do,> she kept saying. <I can't bear it.>

Once she wondered if he would stop if she dared go to him and then she remembered how he had driven her out of the room and thought that perhaps the sight of her might make him worse. Even when she pressed her hands more tightly over her ears she could not keep the awful sounds out. She hated them so and was so terrified by them that suddenly they began to make her angry and she felt as if she should like to fly into a tantrum herself and frighten him as he was frightening her. She was not used to any one's tempers but her own. She took her hands from her ears and sprang up and stamped her foot.

<He ought to be stopped! Somebody ought to make him stop! Somebody ought to beat him!> she cried out.

Just then she heard feet almost running down the corridor and her door opened and the nurse came in. She was not laughing now by any means. She even looked rather pale.

<He's worked himself into hysterics,> she said in a great hurry. <He'll do himself harm. No one can do anything with him. You come and try, like a good child. He likes you.>

<He turned me out of the room this morning,> said Mary, stamping her foot with excitement.

The stamp rather pleased the nurse. The truth was that she had been afraid she might find Mary crying and hiding her head under the bed-clothes.

<That's right,> she said. <You're in the right humor. You go and scold him. Give him something new to think of. Do go, child, as quick as ever you can.>

It was not until afterward that Mary realized that the thing had been funny as well as dreadful—that it was funny that all the grown-up people were so frightened that they came to a little girl just because they guessed she was almost as bad as Colin himself.

She flew along the corridor and the nearer she got to the screams the higher her temper mounted. She felt quite wicked by the time she reached the door. She slapped it open with her hand and ran across the room to the four-posted bed.

<You stop!> she almost shouted. <You stop! I hate you! Everybody hates you! I wish everybody would run out of the house and let you scream yourself to death! You will scream yourself to death in a minute, and I wish you would!>

A nice sympathetic child could neither have thought nor said such things, but it just happened that the shock of hearing them was the best possible thing for this hysterical boy whom no one had ever dared to restrain or contradict.

He had been lying on his face beating his pillow with his hands and he actually almost jumped around, he turned so quickly at the sound of the furious little voice. His face looked dreadful, white and red and swollen, and he was gasping and choking; but savage little Mary did not care an atom.

<If you scream another scream,> she said, <I'll scream too—and I can scream louder than you can and I'll frighten you, I'll frighten you!>

He actually had stopped screaming because she had startled him so. The scream which had been coming almost choked him. The tears were streaming down his face and he shook all over.

<I can't stop!> he gasped and sobbed. <I can't—I can't!>

<You can!> shouted Mary. <Half that ails you is hysterics and temper—just hysterics—hysterics—hysterics!> and she stamped each time she said it.

<I felt the lump—I felt it,> choked out Colin. <I knew I should. I shall have a hunch on my back and then I shall die,> and he began to writhe again and turned on his face and sobbed and wailed but he didn't scream.

<You didn't feel a lump!> contradicted Mary fiercely. <If you did it was only a hysterical lump. Hysterics makes lumps. There's nothing the matter with your horrid back—nothing but hysterics! Turn over and let me look at it!>

She liked the word <hysterics> and felt somehow as if it had an effect on him. He was probably like herself and had never heard it before.

<Nurse,> she commanded, <come here and show me his back this minute!>

The nurse, Mrs~Medlock and Martha had been standing huddled together near the door staring at her, their mouths half open. All three had gasped with fright more than once. The nurse came forward as if she were half afraid. Colin was heaving with great breathless sobs.

<Perhaps he—he won't let me,> she hesitated in a low voice.

Colin heard her, however, and he gasped out between two sobs:

<Sh—show her! She—she'll see then!>

It was a poor thin back to look at when it was bared. Every rib could be counted and every joint of the spine, though Mistress Mary did not count them as she bent over and examined them with a solemn savage little face. She looked so sour and old-fashioned that the nurse turned her head aside to hide the twitching of her mouth. There was just a minute's silence, for even Colin tried to hold his breath while Mary looked up and down his spine, and down and up, as intently as if she had been the great doctor from London.

<There's not a single lump there!> she said at last. <There's not a lump as big as a pin—except backbone lumps, and you can only feel them because you're thin. I've got backbone lumps myself, and they used to stick out as much as yours do, until I began to get fatter, and I am not fat enough yet to hide them. There's not a lump as big as a pin! If you ever say there is again, I shall laugh!>

No one but Colin himself knew what effect those crossly spoken childish words had on him. If he had ever had any one to talk to about his secret terrors—if he had ever dared to let himself ask questions—if he had had childish companions and had not lain on his back in the huge closed house, breathing an atmosphere heavy with the fears of people who were most of them ignorant and tired of him, he would have found out that most of his fright and illness was created by himself. But he had lain and thought of himself and his aches and weariness for hours and days and months and years. And now that an angry unsympathetic little girl insisted obstinately that he was not as ill as he thought he was he actually felt as if she might be speaking the truth.

<I didn't know,> ventured the nurse, <that he thought he had a lump on his spine. His back is weak because he won't try to sit up. I could have told him there was no lump there.>

Colin gulped and turned his face a little to look at her.

<C-could you?> he said pathetically.

<Yes, sir.>

<There!> said Mary, and she gulped too.

Colin turned on his face again and but for his long-drawn broken breaths, which were the dying down of his storm of sobbing, he lay still for a minute, though great tears streamed down his face and wet the pillow. Actually the tears meant that a curious great relief had come to him. Presently he turned and looked at the nurse again and strangely enough he was not like a Rajah at all as he spoke to her.

<Do you think—I could—live to grow up?> he said.

The nurse was neither clever nor soft-hearted but she could repeat some of the London doctor's words.

<You probably will if you will do what you are told to do and not give way to your temper, and stay out a great deal in the fresh air.>

Colin's tantrum had passed and he was weak and worn out with crying and this perhaps made him feel gentle. He put out his hand a little toward Mary, and I am glad to say that, her own tantrum having passed, she was softened too and met him half-way with her hand, so that it was a sort of making up.

<I'll—I'll go out with you, Mary,> he said. <I shan't hate fresh air if we can find—> He remembered just in time to stop himself from saying <if we can find the secret garden> and he ended, <I shall like to go out with you if Dickon will come and push my chair. I do so want to see Dickon and the fox and the crow.>

The nurse remade the tumbled bed and shook and straightened the pillows. Then she made Colin a cup of beef tea and gave a cup to Mary, who really was very glad to get it after her excitement. Mrs~Medlock and Martha gladly slipped away, and after everything was neat and calm and in order the nurse looked as if she would very gladly slip away also. She was a healthy young woman who resented being robbed of her sleep and she yawned quite openly as she looked at Mary, who had pushed her big footstool close to the four-posted bed and was holding Colin's hand.

<You must go back and get your sleep out,> she said. <He'll drop off after a while—if he's not too upset. Then I'll lie down myself in the next room.>

<Would you like me to sing you that song I learned from my Ayah?> Mary whispered to Colin.

His hand pulled hers gently and he turned his tired eyes on her appealingly.

<Oh, yes!> he answered. <It's such a soft song. I shall go to sleep in a minute.>

<I will put him to sleep,> Mary said to the yawning nurse. <You can go if you like.>

<Well,> said the nurse, with an attempt at reluctance. <If he doesn't go to sleep in half an hour you must call me.>

<Very well,> answered Mary.

The nurse was out of the room in a minute and as soon as she was gone Colin pulled Mary's hand again.

<I almost told,> he said; <but I stopped myself in time. I won't talk and I'll go to sleep, but you said you had a whole lot of nice things to tell me. Have you—do you think you have found out anything at all about the way into the secret garden?>

Mary looked at his poor little tired face and swollen eyes and her heart relented.

<Ye-es,> she answered, <I think I have. And if you will go to sleep I will tell you to-morrow.>

His hand quite trembled.

<Oh, Mary!> he said. <Oh, Mary! If I could get into it I think I should live to grow up! Do you suppose that instead of singing the Ayah song—you could just tell me softly as you did that first day what you imagine it looks like inside? I am sure it will make me go to sleep.>

<Yes,> answered Mary. <Shut your eyes.>

He closed his eyes and lay quite still and she held his hand and began to speak very slowly and in a very low voice.

<I think it has been left alone so long—that it has grown all into a lovely tangle. I think the roses have climbed and climbed and climbed until they hang from the branches and walls and creep over the ground—almost like a strange gray mist. Some of them have died but many—are alive and when the summer comes there will be curtains and fountains of roses. I think the ground is full of daffodils and snowdrops and lilies and iris working their way out of the dark. Now the spring has begun—perhaps—perhaps—>

The soft drone of her voice was making him stiller and stiller and she saw it and went on.

<Perhaps they are coming up through the grass—perhaps there are clusters of purple crocuses and gold ones—even now. Perhaps the leaves are beginning to break out and uncurl—and perhaps—the gray is changing and a green gauze veil is creeping—and creeping over—everything. And the birds are coming to look at it—because it is—so safe and still. And perhaps—perhaps—perhaps—> very softly and slowly indeed, <the robin has found a mate—and is building a nest.>

And Colin was asleep.